% SUBMISSION: {"kind": "research_attempt", "category": "economics", "problem_id": "E52-44", "model": {"provider": "OpenAI", "version": "Codex, GPT-6 family", "version_uncertainty": "The serving context exposed the GPT-6 family. Exact serving revisions were not exposed, including for other materially involved Codex agents; no revision is inferred from the model folder.", "reasoning": "Not exposed by the session.", "generated_on": "2026-10-11"}, "other_models": "Separate Codex agents materially constructed or reviewed the equilibrium/filter, policy recursion, examples, original scope, sources and presentation. Their exact serving revisions and reasoning configurations were not exposed. No other material model family is confirmed.", "tools": "File and Git/GitHub tools; primary-source web and PDF inspection; LaTeX compilation and rendered-page inspection when recorded in the production report. Python/Node repository and site checks are publication checks, not proof verification. No Lean or other theorem prover, computer-algebra proof, simulation or executed dynamic-programming implementation is claimed.", "target": "E52-44, Macroeconomic policy and information with dispersed beliefs: under explicit policy, information and competitive-equilibrium-selection restrictions, construct approximating optimal infinite-horizon policies and dynamic programming on an endogenous sufficient state, determine information-welfare effects, compare equally costly public/private disclosure, and characterize transparency changes in stabilization or information allocation.", "scope": "Complete proof claim for the original expressly permitted restricted-class target: a genuine infinite-horizon competitive economy with finitely many signal cohorts, a controlled finite Markov fundamental with strictly positive physical transitions, compact household/producer controls, explicit government finance and information filtrations, and legally committed finite-alphabet disclosure/action prescriptions. The joint calibrated common-information law of the fundamental and private predictive beliefs is derived endogenously, not assumed exogenous. The admissible class includes all public-history-dependent and publicly randomized prescriptions whose only private argument is the current fundamental. The theorem provides sufficient-state closure, continuity, Bellman verification and epsilon-optimal policies, equally costly disclosure comparison, a strict private-disclosure example and an exact transparency/stabilization example. It does not claim a result for every monetary economy, unrestricted private policy memory, informative flexible prices, arbitrary equilibrium selection or universally beneficial information.", "human_contributions": "JiHo Bae requested the problem range, nonduplication, preservation, readable and accurately cited proof, and GitHub/PR submission. No significant human mathematical contribution, substantive proof correction or independent human proof checking is confirmed. OpenAI Codex is the primary mathematical generating attribution; JiHo Bae is credited as submitter and prompting contributor.", "verification": "Separate Codex agents reviewed the equilibrium/filter, dynamic programming, scope, symbolic information examples and primary sources. These are model-assisted reviews. The accompanying final review records identify the complete frozen source, actual audit timing and actual outcomes; this declaration itself does not assert an unperformed or successful final audit. No independent human proof verification, formalization, computer-algebra proof, numerical experiment or executed policy solver is claimed. Compilation, rendering and repository/site checks concern document production and publication software, not mathematical correctness or novelty.", "reproducibility": "Literal Prompt A was applied on 2026-10-11 at 04:17:31 KST before prospective full-proof construction; earlier activity was candidate/model screening, not a completed manuscript. The complete E52-44 original, referenced MA-03/D83-67 economy, actual policy, filled prompt and input hashes were retained. The accompanying reproducibility record identifies the frozen complete proposed source and actual filled inputs, timing and outcome of the final literal Prompt B audit before submission; consult that record for its status. Public source, bibliography, reviews and file hashes accompany the release. Exact serving revisions, reasoning settings and a full original generation transcript are unavailable; no private transcripts or secrets are published."}
% FULL_SOLUTION_SCOPE: {"target": "E52-44, Macroeconomic policy and information with dispersed beliefs: under explicit policy, information and competitive-equilibrium-selection restrictions, construct approximating optimal infinite-horizon policies and dynamic programming on an endogenous sufficient state, determine information-welfare effects, compare equally costly public/private disclosure, and characterize transparency changes in stabilization or information allocation.", "result": "Lemma 1: household and producer optimization, unique normalized competitive price, selected market-clearing allocation, government budget and social-welfare accounting. Theorem 2: calibrated joint common-information state, Bayesian probability law, endogenous filtering and population-belief dynamics, compactness and Feller/reward continuity through null disclosure branches. Theorem 3: Bellman value verification against all admissible public-history randomized commitment policies and uniformly epsilon-optimal stationary Borel policies. Corollary 4: finite approximation, bounded tails and uniform-residual policy certificates. Proposition 5: strict private-disclosure advantage over the optimized public technology at equal cost. Proposition 6: exact posterior stabilization-tax threshold, optimized public transparency and discounted blackout comparison.", "coverage": "The original explicitly permits policy, information and equilibrium-selection restrictions, approximating policies instead of exact attainment, and dynamic programming instead of first-order conditions. The proof covers the declared competitive class: every allowed finite fundamental/cohort configuration, positive physical transition parameter, compact feasible disclosure technology, calibrated initial common law and discount factor in (0,1); channel zeros and null public/private branches; all public-history-dependent and publicly randomized committed prescriptions whose only private argument is the current fundamental; globally optimizing competitive household deviations, normalized prices, government finance and market/resource clearing; endogenous full-history conditional beliefs and their proved sufficient projection; and the original information-welfare, equal-cost public/private and stabilization/allocation questions. The class admits iid physical kernels as a special case, where the exact strict disclosure and stabilization examples use the same equilibrium model and declared control restrictions. Conditional Bellman lemmas or reset examples alone are not the full-solution basis. No result for every macroeconomic economy, unrestricted private-history execution, arbitrary selection, or universally beneficial information is claimed.", "dependencies": "Patrick Billingsley, Weak Convergence of Measures: Applications in Probability, SIAM, 1971, Section 2, printed pp.4-5, for measure determination and metrizability; Section 4, Theorem 4.1, printed p.10, for tightness and relative weak compactness for complete separable metric spaces; applied to probability measures on a compact metric state space, with the model-specific Bayes-calibration subset proved closed. Mankiw and Reis, Imperfect Information and Aggregate Supply, Handbook of Monetary Economics 3A, Chapter 5 (2011), Sections 5.2 and 7.3: motivating policy/information agenda, not the new equilibrium or filtering theorem. Competitive optimization, finite-state Bayes updates, null-branch continuity, discount contraction, epsilon-policy verification and displayed information examples are supplied by proof rather than imported as unverified claims. Rick Durrett, Probability: Theory and Examples, author-posted Version 5, 11 January 2019, Theorems 2.1.22 and 4.1.17: standard Borel targets admit regular conditional distributions, applied to the full-history belief lift.", "human_verification": "None confirmed. Separate Codex-agent checks are model-assisted reviews, not independent human proof verification.", "unverified": "The entire proof has not been independently checked by a human. Lean/formalization, computer-algebra proof, simulation and execution of a dynamic-programming solver were not performed. Symbolic boundary and example checks and theorem-hypothesis reviews are model assisted; their actual source-bound final outcomes are recorded separately. No missing mathematical obligation is asserted within the expressly declared original-target class. The solved/100 percent declaration expresses a full-scope proof claim, not confidence, novelty certification or independent verification.", "novelty": "The earlier E52-44 attempts and their conditional Bellman and simpler information/reset or insurance material are distinguished explicitly. Standard discount contraction and feasible-set inclusion are not claimed novel alone. The new core claim is the derived competitive equilibrium and persistent hidden-state calibrated joint-belief closure, followed by optimization over all legal history policies and the declared information/stabilization results. Primary-source/context and previous-attempt checks were performed. No exhaustive literature/equivalence search or established worldwide novelty claim is made; exact duplicate-clearance evidence is recorded separately."}
% ECONOMICS_PROBLEM: {"id":"E52-44","title":"Macroeconomic policy and information with dispersed beliefs","jel_code":"E52","source_id":"OP-0585"}
% FIRST_POSTED: 2026-10-11
% ATTEMPT_STATUS: solved
% ENTRY_KIND: research_attempt
\documentclass[11pt,a4paper]{article}
\usepackage[T1]{fontenc}
\usepackage[utf8]{inputenc}
\usepackage{lmodern,amsmath,amssymb,amsthm,mathtools,microtype}
\usepackage[margin=24mm]{geometry}
\usepackage{xurl,hyperref,enumitem}
\hypersetup{colorlinks=true,linkcolor=blue,citecolor=blue,urlcolor=blue,
 pdftitle={Policy and disclosure with endogenous dispersed beliefs}}
\setlength{\parindent}{1.1em}
\setlength{\parskip}{3pt}
\setlength{\emergencystretch}{2em}
\setlist{nosep,leftmargin=1.8em}
\newtheorem{theorem}{Theorem}
\newtheorem{lemma}[theorem]{Lemma}
\newtheorem{proposition}[theorem]{Proposition}
\newtheorem{corollary}[theorem]{Corollary}
\newcommand{\R}{\mathbb R}
\newcommand{\E}{\mathbb E}
\newcommand{\Prob}{\mathbb P}
\newcommand{\ind}{\mathbf 1}
\newcommand{\clip}{\operatorname{clip}_{[0,1]}}
\newcommand{\Law}{\mathcal L}
\newcommand{\cP}{\mathcal P}
\newcommand{\cK}{\mathcal K}
\newcommand{\cH}{\mathcal H}
\newcommand{\cF}{\mathcal F}
\newcommand{\dd}{\,d}
\title{Policy and disclosure with endogenous dispersed beliefs}
\author{OpenAI Codex (GPT-6 family)\\
\small Exact serving revision and reasoning setting not exposed\\
\small Submitted by JiHo Bae}
\date{11 October 2026}
\begin{document}
\maketitle

\paragraph{Submission provenance.}
The complete mathematical argument and its contribution disclosures are
retained from the \href{https://github.com/jbaelaw/persistent-private-beliefs-policy-proof/blob/03889470bdc7d04b6b6bbd73c0abfa1e0add4555/paper/Proof.tex}{immutable public source} and
\href{https://github.com/jbaelaw/persistent-private-beliefs-policy-proof/blob/03889470bdc7d04b6b6bbd73c0abfa1e0add4555/paper/Proof.pdf}{proof PDF}. This catalogue wrapper supplies identity and
submission metadata. The full-solution scope is the declared competitive
economy class described in the manuscript. Community review remains pending.

\begin{abstract}
We solve the policy and information problem in E52-44 for a declared class of
competitive economies with a hidden finite fundamental and heterogeneous private
beliefs. Household demands, public accounts and the common posterior transition
are derived from the primitives. A calibrated probability measure is sufficient
for policy, including persistent fundamentals and disclosure channels with zero
entries. The resulting discounted Bellman equation characterizes the value and
stationary Borel policies that approximate it uniformly. Equally costly private
disclosure can strictly outperform the optimized public technology. A separate
example identifies the exact posterior threshold at which transparency changes
the optimal stabilization tax.
\end{abstract}

\section{The question and the declared economy}
The registered problem E52-44 \cite{catalogue} asks, under explicit policy,
information and equilibrium-selection restrictions, for optimal or approximating
policies in an infinite-horizon competitive economy with dispersed information,
dynamic programming on a sufficient state, and welfare comparisons between
equally costly public and private disclosure. Its distribution of individual
states and beliefs must be generated by equilibrium. The accompanying MA-03
entry, D83-67, requires household optimization, Bayes-consistent beliefs and
market clearing; a posterior over current productivity is not sufficient merely
by definition. Mankiw and Reis \cite[Sections 5.2 and 7.3]{mankiwreis} provide the
policy agenda, rather than the theorem proved here.

We specify the restrictions below as permitted by that question. The competitive
price is constant, but the fundamental, private posteriors, aggregate demand and
their joint distribution are endogenous to policy. No assertion about arbitrary
competitive equilibrium correspondences is intended.

There are dates $t=0,1,\ldots$, a discount $0<\beta<1$, and a nonatomic unit
population divided into $g\geq1$ signal groups of masses $w_i>0$, with
$\sum_iw_i=1$. Members of a group share its information; each household remains
an individual price taker. The fundamental belongs to a finite set
$\Theta=\{\theta_1,\ldots,\theta_d\}\subset[1,2]$, with $d\geq1$.
The finite nonempty control set $A$ specifies taxes $\tau_a\geq0$ and transition
matrices $P_a$. Fix $0<c\leq1/d$ and require
\[
 P_a(\ell,\ell')\geq c\qquad(a\in A,\ \ell,\ell'\in[d]).
\]
The next fundamental has this transition conditional on the current fundamental
and control, independently of the remaining past. Iid transitions are allowed;
the theorem also covers persistent ones, for example the two-state matrix
$\left(\begin{smallmatrix}0.9&0.1\\0.1&0.9\end{smallmatrix}\right)$ with $c=0.1$.

There are a numeraire and a produced good. Each household receives $B$ numeraire
units each date, where $B>1+\max_a\tau_a$, and chooses
$(c_i^0,q_i)\in[0,B]\times[0,1]$. There is no storage, saving or individual
physical state carried between dates. Its period utility is
$c_i^0+\theta q_i-q_i^2/2-\kappa Q^2$, where $\kappa\geq0$ is fixed and
$Q=\sum_iw_iq_i$. A competitive firm has the compact production set
$\{(-v,v):-2\leq v\leq2\}$; households own its profits in fixed shares.
The government collects the output tax and consumes exactly
$G=\tau_aQ$ numeraire units. This public consumption has unit social value;
social welfare sums private utilities and this public consumption, counting
the external loss $\kappa Q^2$ once. Individual households cannot influence that
aggregate loss or future policy through their own choices.

\paragraph{Information and commitment.}
Finite nonempty alphabets $Y,Z_1,\ldots,Z_g$ contain a public message and group
messages. A prescription $k$ is a probability array
$k(a,y,z_1,\ldots,z_g\mid\ell)$ for each $\ell$.
The feasible set $\cK$ is any nonempty compact subset of this finite product of
simplexes. Channel entries may be zero. At each date the authority chooses and
publicly announces $k$ using the public history, with optional public
randomization. It observes the current fundamental and executes this announced
array. It may not condition execution on its private past or past private
messages. Households observe the realized $(a,y)$ and their own $z_i$ before
choosing quantities. Thus the tax itself is an observed signal and is included
in the filter. These are explicit commitment and information restrictions on
the authority's adapted policy.

The public history contains previous announcements and public outcomes. A
household also remembers its group's messages. Neither the current fundamental,
realized utility, aggregate orders nor government receipts are separately
observed. The public price is observed. Own choices and numeraire consumption
are functions of information already held and add no signal. No subsequent
unmodeled feedback is supplied. These restrictions prevent a hidden extra
learning channel.
The firm fills aggregate orders after household choices, and the government
administers receipts. Neither observation is shared as household feedback or
used as an additional argument of an admissible prescription.

\section{Competitive allocation and the sufficient common state}
Write $\Delta_c=\{b\in\R^d:b_\ell\geq c,\ \sum_\ell b_\ell=1\}$.
Before current disclosure, $b_i$ denotes group $i$'s posterior over the current
fundamental given its entire information then available. The candidate public
state is the conditional law $\pi$ of $(\ell,b_1,\ldots,b_g)$ on
$S=[d]\times\Delta_c^g$. Its admissible values form
\begin{equation}\label{eq:calibration}
 \Pi=\left\{\pi\in\cP(S):
 \int(\ind_{\{\ell=j\}}-b_{ij})\varphi(b_i)\dd\pi=0
 \quad\forall i,j,\ \forall\varphi\in C(\Delta_c)\right\}.
\end{equation}
This is the Bayes calibration identity: conditional on group $i$'s posterior,
the fundamental has that distribution. Any $\pi_0\in\Pi$ is an admissible
initial condition. It is implemented by drawing $(\ell,b)$ from $\pi_0$ and
showing group $i$ its coordinate $b_i$. The identity for continuous test
functions extends to bounded Borel functions, since continuous functions
determine finite measures on a compact metric space
\cite[Section 2, pp.\ 4--5]{billingsley}. Hence these initial
beliefs are genuine conditional probabilities.

\begin{lemma}[Competitive allocation]\label{lem:equilibrium}
For every admissible information history the normalized competitive price is
$(1/2,1/2)$. If $\widehat b_i$ is the posterior after current disclosure, the
unique household allocation is
\begin{equation}\label{eq:demand}
 q_i=\clip\left(\sum_\ell\theta_\ell\widehat b_{i\ell}-1-\tau_a\right),
 \qquad c_i^0=B-(1+\tau_a)q_i.
\end{equation}
Select firm output $v=Q$ and government consumption $G=\tau_aQ$. This is a
competitive equilibrium, with realized social welfare
\begin{equation}\label{eq:welfare}
 W(\ell,a,\widehat b)=B+(\theta_\ell-1)Q
       -\frac12\sum_iw_iq_i^2-\kappa Q^2.
\end{equation}
\end{lemma}
\begin{proof}
At a normalized price $(p_0,p_1)$, the firm's profit is $(p_1-p_0)v$.
If $p_1>p_0$, its unique supply is $2$, which cannot clear household demand
$Q\in[0,1]$. If $p_1<p_0$, its supply is $-2$, which also cannot clear it.
Thus any equilibrium has $p_0=p_1=1/2$, including exclusion of boundary prices.
Profits are then zero and every $v\in[-2,2]$ is optimal.

At the relative price one, the household budget is
$c_i^0+(1+\tau_a)q_i\leq B$. Utility is increasing in the numeraire, so it
binds. The conditional expected externality is independent of an individual
household's choice. Substitution therefore leaves a strictly concave quadratic in $q_i$ whose
unconstrained maximizer is the argument of the clipping operation in
\eqref{eq:demand}. The bound on $B$ makes every $q_i\in[0,1]$ affordable.
No current choice changes future resources, information or policy; therefore
datewise maximization also maximizes the discounted household objective.
All utilities are bounded, so the discounted sum exists.

The selected $v=Q$ is profit maximizing. Private, production and public
numeraire uses satisfy
$\sum_iw_ic_i^0+Q+G=B$, and produced-good demand equals output. Government
spending equals tax revenue, without a deficit. Adding private utilities and
$G$ gives \eqref{eq:welfare}. The tax
transfer cancels exactly once. Strict concavity gives unique household demands;
the stated firm and government choices complete the selection.
\end{proof}

\begin{theorem}[Endogenous filtering and closure]\label{thm:closure}
The space $\Pi$ is nonempty, compact and metrizable in the weak topology.
For every $\pi_0\in\Pi$ and admissible policy, the competitive economy has
a Bayes-consistent probability law. The conditional common state remains in
$\Pi$, and the formulas below determine its transition, the endogenous
population distribution of allocations and beliefs, and expected welfare.
The transition is jointly Feller in $(\pi,k)$ and the reward is bounded and
jointly continuous, including channels with zero entries.
\end{theorem}
\begin{proof}
For any $b\in\Delta_c$, drawing $\ell$ from $b$ and setting every private
prior equal to $b$ gives a member of $\Pi$. The compact metric space $S$ is
complete and separable. Every probability measure on it is tight, using $S$
itself as the compact set. Billingsley's tightness criterion
\cite[Section 4, Theorem 4.1, p.10]{billingsley} therefore makes $\cP(S)$
weakly compact; the weak topology is metrizable
\cite[Section 2, p.5]{billingsley}. Every integrand in \eqref{eq:calibration} is
continuous, so $\Pi$ is a closed subset and is compact. Taking $\varphi=1$
also shows $\pi(\ell=j)=\int b_{ij}\dd\pi\geq c$.

Write $u=(a,y)$ and $z=(z_1,\ldots,z_g)$. For each group define
\[
 L_i^{k,u,z_i}(\ell)=\sum_{z_{-i}}k(u,z_i,z_{-i}\mid\ell),
 \qquad D_i=\sum_jb_{ij}L_i^{k,u,z_i}(j).
\]
For $D_i>0$ the posterior and next predictive prior are
\begin{equation}\label{eq:filter}
 \widehat b_{i\ell}=\frac{b_{i\ell}L_i^{k,u,z_i}(\ell)}{D_i},
 \qquad b_i'=\widehat b_iP_a.
\end{equation}
Because $b_{ij}\geq c$, a zero denominator means the likelihood is zero at
every fundamental. Every corresponding full profile then has zero probability;
assigning it an arbitrary posterior affects no expectation. In all cases
used with positive probability, $b_i'\in\Delta_c$.

These formulas condition on the full permitted private history. Initially
this follows from calibration and its common-prior implementation. Inductively,
given that history and current $\ell$, the announced prescription has likelihood
$L_i(\ell)$ because execution uses no additional private argument. Bayes' rule
gives the first formula; the observed controlled Markov transition gives the
second. Own choices, spending and the constant price add no signal. This proves
the induction under history-dependent policies and their observed public
randomizations, not only under stationary ones.

The public probability of $u$ is
\[
 h_u(\pi,k)=\int_S\sum_zk(u,z\mid\ell)\dd\pi.
\]
For $h_u>0$, define the next common probability measure $\Phi_u(\pi,k)$ by
\begin{equation}\label{eq:commonfilter}
 \int_S f\dd\Phi_u
 =\frac1{h_u}\int_S\sum_zk(u,z\mid\ell)
       \sum_{j}P_a(\ell,j)
       f\bigl(j,(\widehat b_iP_a)_i\bigr)\dd\pi,
 \qquad f\in C(S).
\end{equation}
This is the conditional law after the finite channel draw and next fundamental
draw. Repeated draws on the initial probability space with fresh independent
random variables construct the economy's law. The state transition is
\begin{equation}\label{eq:kernel}
 \mathsf P(\cdot\mid\pi,k)=\sum_u h_u(\pi,k)\delta_{\Phi_u(\pi,k)}.
\end{equation}
On branches with $h_u=0$ choose any fixed member of $\Pi$; its weight is zero.

To check calibration after a positive public outcome, fix $i,j$ and a
continuous $\varphi(b_i')$. Multiply its desired calibration identity by $h_u$.
Using \eqref{eq:commonfilter} gives
\[
 \int_S\sum_zk(u,z\mid\ell)
 [P_a(\ell,j)-(\widehat b_iP_a)_j]
 \varphi(\widehat b_iP_a)\dd\pi.
\]
Sum over $z_{-i}$ and condition on $b_i$. Calibration replaces the law of
$\ell$ by $b_i$. For every remaining $z_i$, the inner coefficient is
\[
 \sum_\ell b_{i\ell}L_i(\ell)P_a(\ell,j)
       -D_i(\widehat b_iP_a)_j=0.
\]
Thus $\Phi_u\in\Pi$. The lower bound on predictive priors follows directly
from $P_a\geq c$, so no support restriction is lost.

For the realized hidden profile and messages, the actual population measure
at the decision date and its next predictive belief measure are
\begin{equation}\label{eq:mu}
 \mu^{\rm dec}=\sum_iw_i\delta_{(i,\widehat b_i,q_i,c_i^0)},
 \qquad \mu'=\sum_iw_i\delta_{(i,\widehat b_iP_a)}.
\end{equation}
Their public conditional laws are generated by the same draws in
\eqref{eq:commonfilter}; they are not external inputs. The groups are
measurable population intervals, so aggregate quantities are exactly their
weighted sums without an independence or continuum sampling assumption.
For precision, let $H_t$ be the full hidden history of this constructed process.
Its coordinates lie in standard Borel spaces. Regular conditional
distributions therefore exist \cite[Theorems 2.1.22 and 4.1.17]{durrett}.
Take the conditional history law
$\mathfrak b_{it}=\Law(H_t\mid\cF_{it})$ given public history and group $i$'s
initial information and accumulated private signals. The actual full
population measure is
$\mu_t^{\rm full}=\sum_iw_i\delta_{(i,\mathfrak b_{it},q_i,c_i^0)}$.
It is generated by the constructed process; its current-fundamental marginal
is \eqref{eq:filter}, and \eqref{eq:mu} is its decision-belief projection.
Full history beliefs may differ even when $\pi$ agrees. Their remaining
components affect only terms independent of a household's choice, since no
choice changes future opportunities or the aggregate externality. For the
authority, $\pi$ and the legal prescription determine welfare, joint quantities
and the next $\pi$ law. This is the proved sufficient reduction; it does not
claim to reconstruct the whole history-belief measure from $\pi$.

The expected selected-equilibrium reward is explicitly
\begin{equation}\label{eq:reward}
 r(\pi,k)=\int_S\sum_{u,z}k(u,z\mid\ell)
               W(\ell,a,\widehat b(b,k,u,z))\dd\pi.
\end{equation}
It is bounded because $\theta\in[1,2]$ and $q_i,Q\in[0,1]$.

It remains to prove continuity through zero channels. For $f\in C(S)$,
let $\eta_u(\pi,k)$ be the unnormalized measure in
\eqref{eq:commonfilter}. Each term of its test integrand has the form
\[
 k(u,z\mid\ell)\sum_jP_a(\ell,j)
           f\bigl(j,(\widehat b_iP_a)_i\bigr).
\]
If its coefficient has a positive limit along a convergent sequence of
beliefs and channels, every private likelihood used by that profile has
positive limiting mass. Since $b_i\geq c$, all relevant denominators stay
away from zero, so the term converges by ordinary continuity. If the
coefficient tends to zero, the term's absolute value is at most
$\|f\|_\infty k(u,z\mid\ell)$ and tends to zero irrespective of a null
posterior convention. The finite sum is therefore continuous on the compact
product and uniformly continuous there. Weak convergence of $\pi_n$ and
convergence of $k_n$ imply
$\eta_u(\pi_n,k_n)\Rightarrow\eta_u(\pi,k)$ and $h_u(\pi_n,k_n)\to h_u(\pi,k)$.
Where the limit mass is positive, normalization gives continuity of $\Phi_u$.
Where it is zero, for any $v\in C(\Pi)$ the corresponding transition term
is bounded by $\|v\|_\infty h_u(\pi_n,k_n)\to0$. Summing the public
outcomes proves joint Feller continuity of \eqref{eq:kernel}. It also makes
the transition Borel, since the positive-mass branch maps are continuous and
their zero-mass extensions are Borel. Exactly the same weighted argument,
using continuous clipped demands and bounded continuous welfare in place of
$f$, proves joint continuity of \eqref{eq:reward}.
\end{proof}

\section{Optimal policy on the derived state}
All histories and policies in this section obey the information and commitment
restrictions above. Let $\mathsf P(\dd\pi'\mid\pi,k)$ and $r(\pi,k)$ denote
the transition and expected welfare derived in the preceding section. Put
$M=\sup_{\pi,k}|r(\pi,k)|<\infty$ and, for $v\in C(\Pi)$, set
\begin{equation}\label{eq:bellman}
 Q_v(\pi,k)=r(\pi,k)+\beta\int v(\pi')\mathsf P(\dd\pi'\mid\pi,k),
 \qquad Tv(\pi)=\max_{k\in\cK}Q_v(\pi,k).
\end{equation}

\begin{theorem}[Value and approximating policies]\label{thm:bellman}
The Bellman equation $V=TV$ has a unique continuous solution. It equals the
supremum of discounted equilibrium welfare over all admissible public-history
policies, including randomized ones. For every $\varepsilon>0$ there is a
deterministic stationary Borel policy with value at least $V-\varepsilon$
at every $\pi\in\Pi$. If $V_0=0$ and $V_H=T^H0$, then
\begin{equation}\label{eq:tail}
 \|V-V_H\|_\infty\leq\frac{M\beta^H}{1-\beta}.
\end{equation}
For any continuous candidate $v$ with
$\rho=\|Tv-v\|_\infty$, one has
$\|V-v\|_\infty\leq\rho/(1-\beta)$.
\end{theorem}
\begin{proof}
The derived Feller property makes $Q_v$ continuous on the compact product
$\Pi\times\cK$. Its maximum exists, and
$|Tv(\pi)-Tv(\widetilde\pi)|\leq
\sup_k|Q_v(\pi,k)-Q_v(\widetilde\pi,k)|$ proves continuity by uniform
continuity. Probability normalization gives
$\|Tv-Tw\|_\infty\leq\beta\|v-w\|_\infty$.
The iterates from zero satisfy
$\|V_{h+1}-V_h\|_\infty\leq M\beta^h$. They converge uniformly to a
continuous fixed point $V$, with $\|V\|_\infty\leq M/(1-\beta)$ and
\eqref{eq:tail}. Two fixed points must coincide by the contraction inequality.
Applying that inequality to $v$ and $V$ gives
$\|v-V\|_\infty\leq\rho+\beta\|v-V\|_\infty$, proving the residual bound.

For any state and prescription, the fixed point implies
$V(\pi)\geq r(\pi,k)+\beta\E[V(\pi')\mid\pi,k]$.
The equilibrium closure makes this the correct conditional inequality after
every public history. Conditioning also on public randomization and iterating
gives an upper bound on the first $H$ discounted rewards plus
$\beta^H\E V(\pi_H)$. The latter tends uniformly to zero. Thus $V$ bounds
the value of every admissible policy.

For $\eta>0$, compactness and uniform continuity give a finite set
$F=\{k_1,\ldots,k_L\}\subset\cK$ with
$0\leq V(\pi)-\max_{k\in F}Q_V(\pi,k)\leq\eta$ for every $\pi$.
Select the first maximizer in this list. Its selection sets are finite
intersections of strict or weak comparisons of continuous functions, so the
stationary policy $f$ is Borel. Its Bellman deficit
$d(\pi)=V(\pi)-Q_V(\pi,f(\pi))$ lies in $[0,\eta]$.
Conditional telescoping, followed by the bounded-tail limit, yields
\[
 V(\pi)-J_f(\pi)=\E_f\sum_{t\geq0}\beta^td(\pi_t)
 \leq\frac{\eta}{1-\beta}.
\]
Taking $\eta=(1-\beta)\varepsilon$ proves both approximation and equality
with the policy supremum. No exact measurable maximizing selector is needed.
\end{proof}

\begin{corollary}[Finite approximation and a policy certificate]\label{cor:approx}
Fix $H\geq1$. A common finite action set $F\subset\cK$ can be chosen so that
$0\leq TV_h-\max_{k\in F}Q_{V_h}(\cdot,k)\leq\eta$ for
$h=0,\ldots,H-1$. Let $W_0=0$ and
$W_h=\max_{k\in F}Q_{W_{h-1}}(\cdot,k)$. Then
\[
 0\leq V_h-W_h\leq\frac{\eta(1-\beta^h)}{1-\beta}.
\]
Use the first finite maximizer with the appropriate remaining horizon and
then any feasible action after date $H-1$. Its infinite-horizon value $J_H$
satisfies
\begin{equation}\label{eq:finiteerror}
 0\leq V-J_H\leq
 \frac{\eta(1-\beta^H)+2M\beta^H}{1-\beta}.
\end{equation}
More generally, if a finite first-maximizer policy $f$ satisfies
$0\leq Tv-Q_v(\cdot,f(\cdot))\leq\eta$, then
$0\leq V-J_f\leq(2\rho+\eta)/(1-\beta)$.
\end{corollary}
\begin{proof}
There are finitely many continuous functions $Q_{V_h}$, so one sufficiently
fine finite action net works for all of them. Monotonicity gives $V_h\geq W_h$.
If their maximum difference at stage $h-1$ is $e_{h-1}$, insertion of
$V_{h-1}$ between the full and finite maximizations gives
$e_h\leq\eta+\beta e_{h-1}$. Summation proves the first bound.
Finite first-maximizer policies attain $W_H$. The same Bellman upper-bound
argument as before, together with arbitrarily fine action nets, shows that
$V_H$ is the finite-horizon policy supremum. Each discounted tail has absolute
value at most $M\beta^H/(1-\beta)$, giving \eqref{eq:finiteerror}.
For the last claim,
$\|v-Q_v(\cdot,f(\cdot))\|_\infty\leq\rho+\eta$.
Telescoping bounds $\|v-J_f\|_\infty$ by
$(\rho+\eta)/(1-\beta)$; combine with Theorem~\ref{thm:bellman}.
\end{proof}

For a finitely supported initial $\pi$, each finite prescription and finite
message alphabet leave finite support after filtering. A finite action set
and horizon therefore give a finite public decision tree, evaluated by backward
recursion using the explicit filter; branches of probability zero contribute
zero. Compactness guarantees existence of sufficiently fine action nets but
does not supply a computable mesh for an arbitrary compact technology. A
numerical error certificate requires a proved uniform continuity modulus and
validated evaluations. Sampled states alone do not establish a uniform Bellman
residual. The results above are policy and approximation characterizations,
without an unproved complexity claim.

\section{Public and private disclosure at equal cost}
If two feasible technologies satisfy $\cK_1\subset\cK_2$ at the same
resource cost, every policy under the first is available under the second,
with the same equilibrium law. Hence
\begin{equation}\label{eq:inclusion}
 V(\cK_2)\geq V(\cK_1).
\end{equation}
In particular, private disclosure weakly dominates public disclosure when its
feasible arrays include the public arrays by fixing the private message slots.
This is an option comparison, not a claim that forced additional information
always improves welfare. If those embeddings are excluded or costs differ,
\eqref{eq:inclusion} does not establish a ranking.

The following comparison is strict against the optimized public technology.
Take two equal-sized groups, $\theta=1+\omega$ with
$\omega\in\{0,1\}$, initial probability $1/2$ and every row of $P_a$ equal
to $(1/2,1/2)$. Thus each new fundamental is independent of every past history,
and the predictive private priors reset to $1/2$. Set $A=\{0\}$ and $\tau_0=0$.
Fix binary public and private alphabets. The public technology consists of all
public kernels with private messages held constant. The private technology
consists of all joint kernels on these same fixed alphabets. Both are compact
and have zero resource cost, and the latter includes the former. No initial
private information distinguishes the households in the public technology.

\begin{proposition}[A strict private advantage]\label{prop:private}
If $1/2<\kappa<1$, an admissible private experiment has strictly greater welfare
than every public policy. At $\kappa=3/4$ a private experiment yields
\[
 V(\cK_{\rm priv})-V(\cK_{\rm pub})
 \geq\frac{11}{10000(1-\beta)}>0.
\]
\end{proposition}
\begin{proof}
Under public disclosure all households share posterior
$p=\Prob(\omega=1\mid\text{public information})$ and choose $q_i=Q=p$.
Bayes consistency gives $\E p=1/2$ and $\E[\omega p]=\E p^2$. Thus
\[
 \E W_{\rm pub}=B+(1/2-\kappa)\E p^2
 \leq B+\frac{1/2-\kappa}{4}.
\]
The inequality uses $\E p^2\geq(\E p)^2$ and the negative coefficient.
Concealment attains it. This bound holds conditional on every public history
under any admissible public policy, because the current fundamental is iid.
Consequently the optimized discounted public value is
$(B+(1/2-\kappa)/4)/(1-\beta)$. The same upper bound would hold for any
finite public alphabet; no union of varying alphabets is used as a compact
action set here.

Now keep the public message constant and give each group a binary signal,
conditionally independent across groups given $\omega$, with
\[
 \Prob(+\mid\omega=1)=\tfrac12+e,\qquad
 \Prob(+\mid\omega=0)=\tfrac12-e,
 \qquad 0<e<\tfrac12.
\]
Each sign has marginal probability $1/2$. Bayes' rule and household demand
give $q_i=1/2+e$ after plus and $q_i=1/2-e$ after minus. Therefore
\[
 \E q_i^2=\tfrac14+e^2,\qquad
 \E[q_i\mid\omega]=\tfrac12+2e^2(2\omega-1).
\]
Conditional independence yields
$\operatorname{Cov}(q_1,q_2)=4e^4$. Since $Q=(q_1+q_2)/2$,
\[
 \E[\omega Q]=\tfrac14+e^2,
 \qquad \E Q^2=\tfrac14+\tfrac12e^2+2e^4.
\]
Substitution into the actual competitive welfare \eqref{eq:welfare} gives
\[
 \E W_{\rm priv}
 =B+\tfrac18-\tfrac\kappa4
       +\frac{1-\kappa}{2}e^2-2\kappa e^4.
\]
Its gain over the best public experiment is
$\frac{1-\kappa}{2}e^2-2\kappa e^4$, positive whenever
$0<e^2<(1-\kappa)/(4\kappa)$. For $1/2<\kappa<1$ this condition also
ensures $e<1/2$, as required by the channel probabilities.
At $\kappa=3/4$, $e=1/10$, the private payoff is
$B-307/5000$, the best public payoff is $B-1/16$, and their difference is
$11/10000$. Repeating this feasible private experiment each date and summing
the geometric series proves the discounted bound. It is a lower bound on
the optimized private value; optimality of the private witness is unnecessary.
\end{proof}

At fixed tax zero, forced full revelation instead of concealment changes
period welfare by $(1/2-\kappa)/4$. It is strictly harmful when
$\kappa>1/2$, even though enlarging an optional disclosure technology cannot
reduce its optimized value. Information affects both useful individual
responses and the social loss from correlated aggregate demand.

\section{When transparency changes stabilization}
Retain the iid fundamental and initial information of the preceding section,
but set $\kappa=1/4$ and allow taxes $\tau\in\{0,2/5\}$.
Use a fixed binary public alphabet and constant private messages. The feasible
prescriptions have the form
\[
 k(a,y\mid\omega)=L(y\mid\omega)\ind_{\{a=f(y)\}},
\]
where $L$ is any kernel on the fixed alphabet and $f$ is a tax assignment
chosen and announced before the current fundamental is used. There are only
finitely many $f$, so this is a compact technology. The control can use public
history and the generated public message but cannot separately encode the
authority's unannounced private fundamental. Disclosure is again free. A
blackout restricts the message to a constant, making its tax constant
conditional on history. This restriction is essential to the comparison.

\begin{proposition}[The exact tax threshold]\label{prop:tax}
With common posterior $p$, the optimal tax is zero for $0<p<3/5$ and $2/5$
for $p>3/5$. Both tie at $p=3/5$ and at $p=0$. Full public revelation
maximizes discounted welfare over the declared disclosure and tax technology.
Its gain over blackout is
\[
 \frac{41}{400(1-\beta)}.
\]
\end{proposition}
\begin{proof}
Household optimization gives $q=Q=(p-\tau)_+$, with no upper clipping since
$p\leq1$ and taxes are nonnegative. Conditional expected incremental welfare is
\[
 H_\tau(p)=p(p-\tau)_+-\tfrac34(p-\tau)_+^2.
\]
For tax zero this is $p^2/4$. For tax $2/5$ it is zero when $p\leq2/5$,
and when $p>2/5$ it equals
$p^2/4+p/5-3/25$. Comparing the two expressions gives the stated threshold
and all ties. The optimized posterior payoff is therefore
\[
 H(p)=\tfrac14p^2+\max\{0,p/5-3/25\}.
\]
It is convex on $[0,1]$. For any experiment, $\E p=1/2$; the convex chord
bound gives
\[
 \E H(p)\leq\tfrac12H(0)+\tfrac12H(1)=\frac{33}{200}.
\]
Full revelation attains the bound by setting tax zero in the low state and
$2/5$ in the high state. This proves global optimality within the declared
discrete tax technology, rather than only comparison of two candidates.
Under blackout, $p=1/2$ and tax zero is uniquely optimal, giving $H(1/2)=1/16$.
Thus the period gain is $33/200-1/16=41/400$.

All physical transition rows are the same and do not depend on tax. The
next fundamental is independent of current information and controls; every
next predictive prior is $1/2$. Accordingly the same upper bound applies
after every public history, including public randomizations. Repeating the
maximizing prescription each date attains the discounted optimum. The two
values are $(B+33/200)/(1-\beta)$ and $(B+1/16)/(1-\beta)$, proving the
claim. Transparency changes stabilization from no tax under blackout to a
positive tax in the revealed high state.
\end{proof}

\section{Scope, provenance and verification}
\paragraph{Full-solution scope.}
The target is every request in E52-44 under the explicit class and policy
restrictions of Section~1. Lemma~\ref{lem:equilibrium} supplies the selected
competitive equilibrium and accounts; Theorem~\ref{thm:closure} derives its
endogenous state law; Theorem~\ref{thm:bellman} and
Corollary~\ref{cor:approx} characterize infinite-horizon optimal value and
approximating policies. The disclosure results compare equally costly feasible
technologies, give a strict private advantage over optimized public disclosure,
and identify when transparency changes stabilization. The result covers all
finite dimensions, all calibrated initial states, all admissible history
policies, $0<\beta<1$, strictly positive physical transitions and disclosure
arrays with zero cells specified above. It is a solution for this declared
class, as the original permits; it does not resolve every macroeconomic model.

The imported probability facts are measure determination by continuous tests,
metrizability of weak convergence, and weak compactness on a compact metric
space, obtained from the cited parts of Billingsley, and existence of regular
conditional distributions on standard Borel spaces from Durrett. Bayesian
filtering, equilibrium, continuity, Bellman
verification and examples are proved here. The handbook source motivates the
question and is not credited with these results. Earlier repository attempts
contain conditional dynamic programming and simpler information comparisons;
the present proof derives the calibrated hidden-state closure and the
competitive policy problem before applying dynamic programming. No exhaustive
literature or novelty search is claimed.

\paragraph{Contributions and checking.}
OpenAI Codex agents generated the mathematics, source audit and manuscript on
11 October 2026. The platform identifies the GPT-6 family but exposes neither
the exact serving revision nor the reasoning setting of each participating
agent. JiHo Bae specified the rank range, preservation, prose and citation
requirements and requested submission. No significant human mathematical
contribution or independent human proof checking is confirmed. All proofs and
examples remain unverified by an independent human; model-assisted reviews and
document or repository checks do not constitute that verification. No Lean
formalization, simulation or mathematical search computation is used.

The prospective research used the repository's literal Prompt A on
11 October 2026 at 04:17:31 Asia/Seoul, before full-proof construction, with
the complete registered E52-44 and D83-67 inputs and contribution policy.
The initial input packet, proof-development notes and source receipts are
retained with the submission. The accompanying review record identifies the
actual date, frozen source hash and outcome of the required final Prompt B
audit. Exact private transcripts are not
published. Primary Mankiw--Reis and Billingsley texts were consulted for the
specific roles stated above. Boundary checks are included in the proofs;
no independent mathematical certification is claimed.

\section{Completion Estimate}
\noindent\textbf{COMPLETION: 100\%}
For the expressly specified competitive class, the proof establishes equilibrium,
endogenous filtering and recursion, discounted policy approximation, equally
costly public/private comparison and the effect of transparency on stabilization.
These cover the original question's bundled requests under its permitted explicit
restrictions. This is a scope claim, not a confidence or verification percentage.

\begin{thebibliography}{9}
\small\raggedright
\bibitem{catalogue}
\emph{Erdosproblems-llm-hunter}, registered economics statements E52-44,
``Macroeconomic policy and information with dispersed beliefs,'' and D83-67,
``Competitive dynamics with dispersed imperfect information,'' baseline
\texttt{738965af403ff4375a58450fd971c43a817843bc}.
\url{https://github.com/mehmetmars7/Erdosproblems-llm-hunter/tree/738965af403ff4375a58450fd971c43a817843bc/attacks/open_problems/economics/statements}.
\bibitem{mankiwreis}
N. Gregory Mankiw and Ricardo Reis (2011).
``Imperfect Information and Aggregate Supply.''
\emph{Handbook of Monetary Economics} 3A, Chapter 5, 183--229,
Sections 5.2 and 7.3.
\url{https://doi.org/10.1016/S0169-7218(11)03005-X}.
Author-hosted text:
\url{https://personal.lse.ac.uk/reisr/papers/10-HofME.pdf}.
\bibitem{billingsley}
Patrick Billingsley (1971).
\emph{Weak Convergence of Measures: Applications in Probability}.
CBMS-NSF Regional Conference Series in Applied Mathematics 5, SIAM,
Section 2, pp.\ 4--5; Section 4, Theorem 4.1, p.10.
\url{https://doi.org/10.1137/1.9781611970623}.
Consulted text:
\url{https://www.math.ucdavis.edu/~gravner/MAT236A/materials/Billingsley-paper.pdf}.
\bibitem{durrett}
Rick Durrett.
\emph{Probability: Theory and Examples}, author-posted Version 5,
11 January 2019, Theorems 2.1.22 and 4.1.17.
\url{https://services.math.duke.edu/~rtd/PTE/PTE5_011119.pdf}.
\end{thebibliography}
\end{document}
